\documentclass[11pt]{article}
\usepackage[margin=1in]{geometry}
\usepackage{amsmath,amssymb,amsthm,booktabs,array}
\usepackage[T1]{fontenc}
\usepackage{lmodern}
\usepackage[colorlinks=true,urlcolor=blue,citecolor=blue,linkcolor=blue]{hyperref}
\newtheorem{theorem}{Theorem}
\newtheorem{lemma}[theorem]{Lemma}
\newtheorem{proposition}[theorem]{Proposition}
\newtheorem{corollary}[theorem]{Corollary}
\newcommand{\ZC}{Z_C}
\title{Three-leaf counterexamples to Mukherjee and Ray's connected-forcing crossing question}
\author{Haobo Ma\\\small ChronoAI Pte Ltd, Singapore
\and Wenlin Zhang\\\small National University of Singapore\\\small ChronoAI Pte Ltd, Singapore}
\date{20 November 2026}
\begin{document}
\maketitle
\begin{abstract}
\begin{sloppypar}
Using Brimkov's earlier starlike enumeration (May 2017, Proposition 7.11),
we answer Mukherjee and Ray's Problem 5.3 in \emph{On the connected zero forcing polynomial of finite
trees}, arXiv:2611.10985v1 (2026), Section 5, p.~21, negatively even for
two three-leaf trees and even when their order may be exchanged.
For every $n\geq10$ we construct such a pair: its coefficient differences
at sizes $4,6,n-2$ are $-2,1,-1$, so neither order admits the requested
single crossing. The construction follows from an explicit enumeration
of connected forcing sets in a three-armed tree. We also explain the
smaller, ordered path--star obstruction from the posers' own examples.
For the neighboring unicyclic problem, a complete five-vertex classification
shows that a fixed-leaf class can have neither a least nor a greatest polynomial.
\end{sloppypar}
\end{abstract}

\section{Failure in both orders with three leaves}
Mukherjee and Ray ask whether the connected zero forcing coefficient
sequences of two equally large trees always cross in a prescribed
direction~\cite[Problem 5.3]{MR}. The answer remains negative if the trees
must have the same number of leaves and either order may be chosen.
Three-armed trees already give an obstruction at every order at least ten.

All graphs are finite, simple and undirected. A colored vertex with exactly
one uncolored neighbor forces that neighbor to become colored. A set $S$
is a \emph{connected zero forcing set} if its successive forces color the
whole graph and the initially induced subgraph $G[S]$ is connected.
This is the convention in~\cite{BD}; singletons are connected.
Write $\ZC(G,i)$ for the number of these sets of size $i$ and
$\ZC^G(x)=\sum_{i=1}^{|V(G)|}\ZC(G,i)x^i$.
For two trees $T_1,T_2$ on $n$ vertices, the requested cutoff is an integer
$\ell$ with $1\leq\ell\leq n$ such that
\begin{align}
\ZC(T_1,i)&\leq\ZC(T_2,i) &&(1\leq i\leq\ell),\label{eq:low}\\
\ZC(T_1,i)&\geq\ZC(T_2,i) &&(\ell<i\leq n).\label{eq:high}
\end{align}
The source allows different leaf counts; equal counts are also admissible.
The possibility of exchanging the order is a strengthening considered here,
not a separate assertion made by the source.

Let $S_{h_1,\ldots,h_k}$ be the tree consisting of a center and $k$ arms,
each a path of $h_j$ edges, with distinct arms meeting only at the center.
Put $F_h(x)=1+x+\cdots+x^{h-1}$ for $h\geq1$.

\begin{theorem}\label{thm:main}
Let $a,b\geq4$ be integers, $n=a+b+2$, and
\[
A=S_{1,a,b},\qquad B=S_{2,2,a+b-3}.
\]
Both trees have $n$ vertices and three leaves. If
$d_i=\ZC(A,i)-\ZC(B,i)$, then
\[
d_4=-2,\qquad d_6=1,\qquad d_{n-2}=-1.
\]
Neither $(A,B)$ nor $(B,A)$ has a cutoff satisfying
\eqref{eq:low}--\eqref{eq:high}.
\end{theorem}

We use Brimkov's enumeration of spiders
\cite[Proposition 7.11, pp.~168--169]{B}, later expressed polynomially by
Mukherjee and Ray~\cite[Lemma 3.3, p.~5]{MR}.
We include the elementary counting proof.
\begin{lemma}\label{lem:enumeration}
For $k\geq3$ and positive arm lengths $h_1,\ldots,h_k$,
\[
\ZC^{S_{h_1,\ldots,h_k}}(x)
=x^k\left(x\prod_{j=1}^k F_{h_j}(x)
+\sum_{j=1}^k\prod_{t\ne j}F_{h_t}(x)\right).
\]
\end{lemma}
\begin{proof}
A connected set avoiding the center lies in one arm. Even if it colors
that entire arm and then the center, the center still has at least two
uncolored arms and cannot force further. Thus every connected forcing set
contains the center. Its intersection with each arm is an initial segment
of length $r_j$, where $0\leq r_j\leq h_j$.
If two or more segments are empty, the center can never start either
empty arm. If at most one is empty, each nonempty segment extends along
its arm; the center then starts the one remaining arm, if present.
These conditions are necessary and sufficient.

Write $L_h=x+\cdots+x^h=xF_h$. Counting the center and either no empty
arm or one specified empty arm gives
$x(\prod_j L_{h_j}+\sum_j\prod_{t\ne j}L_{h_t})$.
Factoring the powers of $x$ yields the formula.
\end{proof}

\begin{proof}[Proof of Theorem~\ref{thm:main}]
The vertex counts are $1+1+a+b=1+2+2+(a+b-3)=n$.
Lemma~\ref{lem:enumeration} gives
\begin{align}
x^{-3}\ZC^A(x)&=F_a+F_b+(1+x)F_aF_b,\label{eq:A}\\
x^{-3}\ZC^B(x)&=(1+x)^2+(2+3x+2x^2+x^3)F_{a+b-3}.\label{eq:B}
\end{align}
Since $a,b\geq4$, the coefficients of $F_aF_b$ at degrees $0,1,2,3$
are $1,2,3,4$. Thus the coefficients at degree $1$ in
\eqref{eq:A} and \eqref{eq:B} are respectively $2+2+1=5$ and
$2+2+3=7$; those at degree $3$ are $2+4+3=9$ and $2+3+2+1=8$.
Here $F_{a+b-3}$ has coefficient one through degree $3$.
These calculations give $d_4=-2$ and $d_6=1$.

At degree $a+b-3$, the terms $F_a+F_b$ contribute zero.
The two highest diagonals needed from $F_aF_b$ have coefficients
$[x^{a+b-3}]F_aF_b=2$ and $[x^{a+b-4}]F_aF_b=3$.
Consequently \eqref{eq:A} has coefficient $5$ there.
The degree of $F_{a+b-3}$ is $a+b-4$, so in \eqref{eq:B} exactly
the multiplier terms of degrees $1,2,3$ contribute, giving $3+2+1=6$.
This proves $d_{n-2}=-1$.

For $(A,B)$, $d_6>0$ requires $\ell<6$, whereas $d_{n-2}<0$
requires $\ell\geq n-2\geq8$, a contradiction.
For $(B,A)$, $-d_4>0$ requires $\ell<4$, whereas $-d_6<0$
requires $\ell\geq6$. This is again impossible.
\end{proof}

\begin{corollary}\label{cor:allorders}
For every $n\geq10$, two $n$-vertex trees with three leaves fail the
crossing condition in both orders. At $n=10$ one may take
$A=S_{1,4,4}$ and $B=S_{2,2,5}$, with
\begin{align*}
\ZC^A(x)&=3x^3+5x^4+7x^5+9x^6+7x^7+5x^8+3x^9+x^{10},\\
\ZC^B(x)&=3x^3+7x^4+8x^5+8x^6+8x^7+6x^8+3x^9+x^{10}.
\end{align*}
\end{corollary}
\begin{proof}
Choose $a=4$ and $b=n-6$. The two displayed polynomials follow by
substituting $a=b=4$ in \eqref{eq:A}--\eqref{eq:B}.
\end{proof}
The existence of a greatest three-leaf tree, proved in
\cite[Corollary 4.12]{MR}, is compatible with this result:
a common coefficientwise upper bound need not impose a single crossing
on every pair below it.

\section{The small ordered obstruction and elementary families}
\begin{corollary}\label{prop:families}
For the path $P_n$ ($n\geq2$), star $K_{1,n-1}$ ($n\geq4$), and
cycle $C_n$ ($n\geq3$), respectively,
\begin{align}
\ZC^{P_n}(x)&=2x+\sum_{i=2}^n(n-i+1)x^i,\label{eq:path}\\
\ZC^{K_{1,n-1}}(x)&=(n-1)x^{n-1}+x^n,\label{eq:star}\\
\ZC^{C_n}(x)&=n\sum_{i=2}^{n-1}x^i+x^n.\label{eq:cycle}
\end{align}
For $p\geq3$, $n\geq p+1$, the starlike tree with $p-1$ unit arms
and one arm of length $n-p$ has polynomial
\begin{equation}
\ZC^{S_{1,\ldots,1,n-p}}(x)=p\sum_{i=p}^{n-1}x^i+x^n.\label{eq:broom}
\end{equation}
More generally, $S_{h_1,\ldots,h_k}$, $k\geq3$, has connected forcing
number $k$ and exactly $k$ minimum connected forcing sets.
\end{corollary}
The path and cycle formulas are Brimkov's Propositions 7.7 and 7.8~\cite{B};
the star is a specialization of his Proposition 7.11.
Equation~\eqref{eq:broom} is also Mukherjee and Ray's Corollary 3.8~\cite{MR}.
The minimum-size count is already given by Brimkov and Davila's
Theorem 13 and Corollary 14~\cite{BD}.
\begin{proof}
Connected subsets of a path are intervals. Every interval of size at least
two forces from its boundaries; there are $n-i+1$ intervals of size $i$.
Of the singleton sets, exactly the two endpoints force.

In a star, a connected set of size at least two contains the center.
It forces exactly when it omits at most one leaf. There are $n-1$ proper
such sets, all of size $n-1$. A singleton cannot force: the center
has at least three uncolored neighbors, and a leaf can color the center
but leaves at least two other leaves uncolored. The smaller stars
$K_{1,1}=P_2$ and $K_{1,2}=P_3$ have the path formulas.

In a cycle, a singleton cannot force. A proper connected subset with at
least two vertices is a cyclic interval and forces from its boundaries.
For each size $2\leq i\leq n-1$ there are $n$ such intervals.
The full set contributes $x^n$ in each case.

For \eqref{eq:broom}, put $h=n-p$ in Lemma~\ref{lem:enumeration}.
It gives $x^p((x+p-1)F_h+1)$, whose coefficients are $p$ at
degrees $p,\ldots,n-1$ and one at degree $n$.
For general positive arm lengths, the same lemma has smallest exponent
$k$ and coefficient $k$, since each $F_{h_j}$ has constant term one.
\end{proof}

\begin{corollary}\label{cor:small}
The smallest order of a counterexample to the literal ordered question
is four. For every $n\geq4$, $(P_n,K_{1,n-1})$ has no cutoff, whereas
the reversed pair has the unique cutoff $n-2$.
\end{corollary}
\begin{proof}
The path-first singleton comparison would be $2\leq0$ for every
permitted cutoff. In the reverse order the star coefficients are zero
and the path coefficients positive through index $n-2$; at index
$n-1$ they are $n-1>2$, and at $n$ they agree. Hence the unique
reverse cutoff is $n-2$. At orders one, two and three there is just
one tree up to isomorphism, so every pair has equal coefficient sequences.
\end{proof}
The two four-vertex polynomials are already printed by Mukherjee and
Ray in Section 4, p.~7~\cite{MR}. Their example supplies the literal
ordered obstruction. Theorem~\ref{thm:main} removes both dependence on
the order and the difference in leaf counts.

\section{The unicyclic extremal question}
Problem 5.2~\cite[p.~21]{MR} asks both for unicyclic enumeration and for
coefficientwise least and greatest graphs at fixed order and leaf count.
The latter question has a negative answer already in a two-element class.

\begin{corollary}\label{cor:unicyclic}
Among connected simple unicyclic graphs with five vertices and two leaves,
there is neither a coefficientwise least nor a coefficientwise greatest
connected zero forcing polynomial.
\end{corollary}
The triangle with leaves at different vertices also appears in
Brimkov's discussion of greedoids~\cite[Section 7.2, pp.~159--160]{B}.
\begin{proof}
The unique cycle has length three: a cycle of length four or five leaves
too few vertices for two leaves. The other two vertices must both be
leaves directly attached to the triangle, since a pendant path of length
two would give just one leaf. Up to isomorphism there are exactly two
graphs: the two leaves attach to the same triangle vertex or to different
triangle vertices.

Label the triangle $0,1,2$ and the leaves $3,4$.
Let $U$ have additional edges $03,04$, and let $V$ have edges $03,14$.
Neither graph has a forcing singleton. No edge-pair of $U$ forces:
coloring along a leaf edge stalls at $0$, while a triangle edge can
color the triangle but leaves both leaves uncolored. The forcing pairs
of $V$ are exactly $\{0,2\}$ and $\{1,2\}$; the other three edges
stall with two uncolored neighbors at each possible forcing vertex.
For triples the complete lists are
\begin{align*}
U:&\quad \{0,1,3\},\{0,1,4\},\{0,2,3\},\{0,2,4\},\\
V:&\quad \{0,1,2\},\{0,1,3\},\{0,1,4\},\{0,2,3\},\{1,2,4\}.
\end{align*}
For $U$, the other connected triples are the triangle and $\{0,3,4\}$;
both stall with two uncolored neighbors at $0$. For $V$, the displayed
list contains every connected triple, and each forces by filling the
triangle and then its remaining leaves.

A connected four-vertex subset forces the remaining vertex. There are
four in $U$, obtained by omitting $1,2,3$ or $4$, and three in $V$,
obtained by omitting $2,3$ or $4$. The full set always qualifies. Thus
\[
\ZC^U(x)=4x^3+4x^4+x^5,\qquad
\ZC^V(x)=2x^2+5x^3+3x^4+x^5.
\]
At size two $V$ is larger; at size four $U$ is larger. Since these are
the only isomorphism classes, no third member can be a common least
or greatest. This proves the assertion about the entire class.
\end{proof}

\section{Source coverage and verification scope}
The following table separates the source's questions and prior results
from the consequences proved above.
\begin{center}
\small
\begin{tabular}{@{}p{0.28\linewidth}p{0.65\linewidth}@{}}
\toprule
Source item in~\cite{MR} & Coverage\\
\midrule
Lemma 3.3, p.~5 & Reproved; prior proof Brimkov, 2017, Proposition 7.11;
later polynomial form Mukherjee--Ray, 2026.\\
Corollaries 3.6, 3.8 & Starlike minimum count and one-long-arm formula
reproved in Corollary~\ref{prop:families}; prior proofs
Brimkov--Davila, 2016, and Mukherjee--Ray, 2026, respectively.\\
Section 4 examples, p.~7 & Path--star polynomials reproved;
prior formulas Brimkov, 2017; later examples Mukherjee--Ray, 2026.\\
Corollary 4.12, p.~15 & Prior proof cited; greatest three-leaf tree.
Its existence does not prevent Theorem~\ref{thm:main}.\\
Conjecture 5.1, p.~20 & Not covered: eventual absence of a greatest
tree with a fixed leaf count $p\geq5$.\\
Problem 5.2, p.~21 & Extremal clause answered: no, by
Corollary~\ref{cor:unicyclic}. Cycle formula reproved; prior proof
Brimkov, 2017, Proposition 7.8. General enumeration not covered.\\
Problem 5.3, p.~21 & Answered: no. Theorem~\ref{thm:main} gives
equal-leaf counterexamples in both orders; Corollary~\ref{cor:small}
gives the smallest literal ordered example.\\
Order exchange variant & Not posed separately in the source;
obstruction proved by Theorem~\ref{thm:main}.\\
Problem 5.4, p.~21 & Partial coverage: the displayed tree and cycle
families and comparisons; arbitrary-graph enumeration not covered.\\
\bottomrule
\end{tabular}
\end{center}
Newton and Brimkov's non-unimodality theorem for connected forcing
polynomials on general graphs~\cite[Corollary 2]{NB} concerns one
coefficient sequence, rather than crossings between two tree sequences.
Their Questions 1--3 concern ordinary zero forcing polynomials and are
not settled by these connected-forcing counts.

\paragraph{Mathematical verification.}
The declaration \texttt{result : $\neg$ claim} in the
\href{https://github.com/the-omega-institute/trureturing/blob/a165df4cad5a5538260a4e40d6b865a3db48795d/D5/S3/Combinatorics/Graph/ZeroForcing/ConnectedForcingCrossingRefutation.lean}{pinned Lean source}
negates the full ordered assertion for arbitrary finite vertex types,
using the four-vertex path and star. It covers the literal negative answer
in Corollary~\ref{cor:small}. The imported
\href{https://github.com/the-omega-institute/trureturing/blob/a165df4cad5a5538260a4e40d6b865a3db48795d/D5/S3/Combinatorics/GeneralizedPetersen/ZeroForcingThree.lean}{\texttt{IsZeroForcing}}
defines ordinary forcing via \texttt{Black.init} and \texttt{Black.force}.
The three-leaf and unicyclic obstructions, the enumeration,
all-order formulas, reverse cutoff, and minimality have the written proofs
above. The accompanying \texttt{verify.py} independently enumerates every
vertex subset of the displayed three-leaf pairs at orders $10$ through
$14$, paths and stars at orders $2$ through $10$, and cycles at orders
$3$ through $10$. It also checks all labeled trees of orders $1$ through
$4$ and the polynomial identities for every $4\leq a\leq b\leq30$.
It exhausts all $1{,}024$ labeled simple graphs of order five, identifies
all $90$ unicyclic two-leaf graphs and their two isomorphism classes,
and tests all $2{,}880$ initial subsets in that class.
These are finite checks, separate from the all-parameter proofs.

\begin{thebibliography}{9}\small
\bibitem{MR} S. Mukherjee and P. Ray,
\emph{On the connected zero forcing polynomial of finite trees},
arXiv:2611.10985v1 (7 November 2026), Lemma 3.3 and Sections 4--5.
\href{https://arxiv.org/abs/2611.10985v1}{arXiv record}.
\bibitem{B} B. Brimkov, \emph{Graph Coloring, Zero Forcing, and Related
Problems}, Ph.D. thesis, Rice University (May 2017), Chapter 7,
Propositions 7.7, 7.8 and 7.11, pp.~165--169.
\href{https://repository.rice.edu/server/api/core/bitstreams/c29c77d5-534a-4c5d-bd83-09121784af53/content}{Rice University text}.
\bibitem{BD} B. Brimkov and R. Davila,
\emph{Characterizations of the Connected Forcing Number of a Graph},
arXiv:1604.00740v1 (4 April 2016), Section 1, Theorem 13 and Corollary 14.
\href{https://arxiv.org/abs/1604.00740v1}{arXiv record}.
\bibitem{NB} T. Newton and B. Brimkov,
\emph{Zero forcing is not unimodal}, arXiv:2611.11352v1
(8 November 2026), Corollary 2 and Section 3.
\href{https://arxiv.org/abs/2611.11352v1}{arXiv record}.
\end{thebibliography}
\smallskip
\noindent\small Copyright 2026 Haobo Ma and Wenlin Zhang.
Licensed under Creative Commons Attribution 4.0 International (CC BY 4.0).
\end{document}
